\documentclass[11pt]{article}

% Change "review" to "final" to generate the final (camera-ready) version.
\usepackage[final]{acl}

% Standard package includes
\usepackage{times}
\usepackage{latexsym}
\usepackage[T1]{fontenc}
\usepackage[utf8]{inputenc}
\usepackage{microtype}
\IfFileExists{inconsolata.sty}{\usepackage{inconsolata}}{}
\usepackage{graphicx}
\usepackage{booktabs}

\title{Diagnose Before Simplifying: Operation-Aware Biomedical Text Simplification}

\author{Sruthilaya \and Rishabh \and Sophakotra \and Zihao \\
  NLP Course Research Project \\
  \texttt{\{sruthilaya, rishabh, sophakotra, zihao\}@northeastern.edu}}

\begin{document}
\maketitle

\begin{abstract}
Biomedical abstracts are written for clinicians, not patients.
Simplification systems that treat every span the same way either
under-edit or over-edit: rule-based lexical substitution is safe but
structurally limited to swapping terms, while unguided LLM prompting can
rewrite freely but risks dropping warnings and dosages in the process. We
reframe simplification as a diagnostic task: before rewriting a span,
classify which operation it requires-Substitution, Explanation, or
Generalization-and route it accordingly. On PLABA, pseudo-labeled
operation distributions show that 56.9\% of required edits are Explanation
or Generalization, operations a purely lexical baseline cannot perform by
construction, motivating generation beyond simple substitution. We present
the first stage of this pipeline: five handcrafted complexity detectors, a
pseudo-labeling scheme deriving span-level operation labels from aligned
source-target pairs, a TF-IDF classifier baseline, and safety-preservation
metrics measuring what SARI misses. On a random 50-sentence sample,
unguided LLM prompting achieves the highest SARI (22.28 vs.\ 16.82 for no
simplification) at near-neutral compression (0.945), with all baselines
preserving warnings perfectly; however, a warning-stratified stress test
(n=30 warning-bearing sentences) reveals unguided prompting is the only
baseline that drops safety-critical content (0.950 preservation vs.\
1.000 for rule-based substitution and no simplification)- demonstrating
that random-sample evaluation alone is insufficient to detect the safety
failures that motivate our operation-aware design.
\end{abstract}

\section{Introduction}

Biomedical abstracts are written for clinicians, not patients. They are dense
with medical jargon, long and syntactically complex sentences, numerical risk
data, and safety warnings that a lay reader cannot follow. This creates a
health literacy gap: the people who most need to understand a treatment, a
risk, or a warning are often unable to read the text that describes it. Low
health literacy is linked to worse health outcomes \citep{berkman2011health},
so making biomedical text readable is a practical clinical concern, not only
a linguistic one.

Most automatic simplification systems treat every kind of complexity the same
way. A sentence is fed into a model and a simpler sentence is returned, whether
the difficulty came from a single jargon term, a tangled clause structure, or a
critical dosage. Directly prompting a language model improves surface readability, but prior
work and our own preliminary evaluation show it can drop medical warnings,
weaken risk language, and remove clinically important numbers---failures that are easy to miss under naive random-sample evaluation (Section~5). In other words, uniform simplification
buys readability at the cost of safety.

We take a diagnostic approach: classify first, then simplify. Before any text is
rewritten, we identify what type of complexity a span carries, and only then
apply the operation suited to it. A jargon term can be replaced from a
controlled vocabulary with no model involved. An unfamiliar concept can be
explained. An overly specific statement can be generalized. Routing each span to
the right operation lets us simplify aggressively where it is safe and stay
conservative where information must be preserved.

Our central research question is: \emph{does predicting the required
simplification operation at the span level reduce safety-critical information
loss, compared to rule-based, trained model, and direct LLM simplification
baselines?} We measure this with custom preservation metrics for entities,
warning phrases, and numerical expressions, alongside standard readability and
similarity measures.

\section{Data}

We use PLABA, the Plain Language Adaptation of Biomedical Abstracts dataset
\citep{attal2023plaba}, the official benchmark at the TREC 2023 and 2024 PLABA
tracks. PLABA contains 750 PubMed abstracts retrieved to answer consumer health
questions, each manually rewritten into plain language by at least one
biomedical expert. We work with a train, validation, and test split of 635,
138, and 148 abstract-level source-target pairs. For sentence-level analysis
(pseudo-labeling, complexity detection, and classifier training), we further
align PLABA's underlying \texttt{data.json} on shared sentence indices between
source abstracts and their adaptations, yielding 6{,}307 leakage-free
sentence-level training pairs (Section 5).

An exploratory analysis of the 635 training abstracts shows the
complexity is multidimensional: $87.9\%$ contain medical jargon,
$82.7\%$ are syntactically complex, $60.9\%$ carry numerical risk data,
and $33.1\%$ contain safety-critical warnings, with $60.8\%$ of abstracts
exhibiting three or more of these simultaneously
(Table~\ref{tab:cooccurrence}). A system applying one operation to every
span will therefore mishandle most inputs, motivating our operation-aware
design.

PLABA does not ship with the span-level operation labels our classifier needs,
so we derive them through pseudo-labeling from the aligned source and reference
sentences (Section~\ref{sec:classifier}). We scope the labels to three of
PLABA's five original operations \citep{attal2023plaba}---Substitution,
Explanation, and Generalization---which together account for over 90\% of
observed replacements \citep{ondov2025lessons} (Section~\ref{sec:classifier}
gives the full breakdown). The Consumer Health Vocabulary (CHV) and UMLS are
used as external knowledge bases in the pipeline rather than as training data.

\begin{table}[t]
  \centering
  \begin{tabular}{crr}
    \toprule
    \textbf{Detectors firing} & \textbf{Abstracts} & \textbf{\%} \\
    \midrule
    0 & 9 & 1.4 \\
    1 & 67 & 10.6 \\
    2 & 173 & 27.2 \\
    3 & 277 & 43.6 \\
    4 & 109 & 17.2 \\
    \bottomrule
  \end{tabular}
  \caption{Complexity co-occurrence on the PLABA training set ($n=635$
    abstracts). Most abstracts carry three or more simultaneous
    complexity types.}
  \label{tab:cooccurrence}
\end{table}

\section{Complexity Detection}
\label{sec:complexity}

Stage 1 of the pipeline is pure classical NLP: no learned model and no
prompting. The design principle is that complexity in biomedical text is not one
problem but several overlapping ones, and each signal is best captured by the
paradigm suited to it. We implement five handcrafted detectors, each targeting a
distinct complexity type.

\paragraph{Numerical expressions.} A regular-expression extractor flags
dosages, risk ratios, confidence intervals, $p$-values, and percentages.
Numbers are exactly what uniform simplification tends to drop or distort,
and a dropped dosage is a safety failure, not a stylistic one.

\paragraph{Warning cues.} A handbuilt lexicon detects safety-critical
language across six categories (contraindications, risk alerts,
avoidance, interactions, overdose, adverse effects), giving full
auditability and zero hallucination risk. This lexicon also grounds our
warning-preservation metric (Section~\ref{sec:metrics}).

\paragraph{Biomedical named entities.} We use SciSpaCy's
\texttt{en\_ner\_bc5cdr\_md} \citep{neumann2019scispacy} to detect diseases
and chemicals; general-domain models miss terms like \emph{nephrotoxicity}.
Detected entities are the primary targets for Substitution.

\paragraph{Syntactic depth.} A dependency-tree depth calculator flags
structurally complex sentences, since parse-tree depth captures genuine
nesting that sentence length alone penalizes unfairly.

\paragraph{UMLS jargon.} A matcher flags medical terms absent from common
English. For this preliminary stage we use a morphological fallback
(biomedical suffixes and drug-naming conventions); the full QuickUMLS
interface is implemented and ready for integration.

No single paradigm captures everything: regular expressions catch numbers but
miss entities, lexicons catch warnings but miss structure, and NER catches
entities but misses jargon that is not a named entity. Using five detectors is a
deliberate design choice reflecting the multidimensional nature of the problem
documented in Table~\ref{tab:cooccurrence}. Because these detectors rely on
handcrafted patterns, their flag rates are a conservative lower bound on the true
complexity present.

\section{Operation Classification}
\label{sec:classifier}

Predicting the operation a span requires is the research centerpiece. Crucially,
the classifier sees only the \emph{source} sentence: at inference time no target
exists, so any feature derived from the reference would be leakage.

\paragraph{Pseudo-labeling.} Because PLABA lacks public span-level operation
labels, we derive them by weak supervision from aligned source--target pairs. We
sentence-align \texttt{data.json} on shared sentence indices, iterate every
available adaptation to maximize volume, and exclude all validation and test
PMIDs to prevent leakage. Each pair is labeled by a length-ratio rule with an
edit-distance tiebreak: a target at least $15\%$ longer is an Explanation, one at
least $15\%$ shorter is a Generalization, and a similar-length pair with
sufficient character-level change is a Substitution. This yields $6{,}307$
labeled training pairs ($43\%$ Substitution, $34\%$ Explanation, $23\%$
Generalization). Deriving labels heuristically from parallel text is a standard
weakly supervised approach when gold labels are inaccessible.

\paragraph{TF-IDF + Logistic Regression.} The first rung of a deliberate
progression (TF-IDF $\rightarrow$ BERT-base $\rightarrow$ BioBERT) uses
$(1,2)$-gram TF-IDF features (\texttt{max\_features}${=}50{,}000$) with a
class-weighted logistic regression to handle the minority Generalization class.
Understanding what this model can and cannot represent is instructive: rare
jargon tokens receive high IDF weight, which helps flag Substitution; bigrams
like ``which is'' signal Explanation; but Generalization is defined by removed
specificity---an \emph{absence}---which bag-of-words features cannot represent.
This is precisely why we expect contextual models to help.

\paragraph{Comparison to the literature.} Our three-operation scope
follows PLABA's original five-operation taxonomy \citep{attal2023plaba},
covering the 90\% of replacements that are Substitution, Explanation, or
Generalization rather than the rarer Omission/Exemplification categories
\citep{ondov2025lessons}. Operation classification is itself known to be
hard---\citet{ondov2025lessons} report human inter-annotator agreement of
only 0.46 F1 and baseline classifier scores of $0.33$--$0.39$ macro-F1---
against which our pseudo-label-trained TF-IDF baseline ($0.424$ macro-F1)
compares favorably.

\section{Baselines and Preliminary Results}
\label{sec:metrics}

We evaluate a ladder of simplification baselines and, for the safety
dimension, implement three preservation metrics from scratch (entity,
warning, and numerical expression preservation), reusing the detectors
from Section~\ref{sec:complexity}.

\paragraph{Baseline ladder.} Baseline~1 performs no simplification and
serves as the safety ceiling: it preserves all warnings by definition.
Baseline~2 applies rule-based CHV substitution with sentence
splitting---pure NLP with zero hallucination risk, but with grammatical
artifacts from context-unaware substitution (e.g., ``is should not be
used''). Baseline~3 is a direct, unguided LLM prompt (FLAN-T5-base),
representing the ``simplify everything uniformly'' approach our
operation-aware system is designed to improve on. We run it locally on
open weights to keep the baseline fully reproducible.

\paragraph{Preliminary results.} Table~\ref{tab:prelim_results} reports
SARI, FKGL, compression ratio, and warning preservation on a random
50-sentence sample of the PLABA validation set. Unguided LLM prompting
achieves the highest SARI (22.28 vs.\ 16.82 for no simplification) at
near-neutral compression (0.945), and all three baselines preserve
warnings perfectly on this random sample---an artifact of sample size,
since only $33.1\%$ of sentences contain warning language
(Section~\ref{sec:complexity}) and a 50-sentence random draw may by
chance contain few or no sentences where a warning is actually dropped.

To test the safety claim directly, we construct a targeted stress test
restricted to warning-bearing sentences only ($n=30$,
Table~\ref{tab:stress_test}). Here the pattern changes: unguided LLM
prompting is the only baseline that drops safety-critical content
(0.950 preservation, versus 1.000 for both no-simplification and
rule-based substitution). This discrepancy between the random-sample and
stress-test results is itself a finding: random-sample evaluation on
rare-but-critical content is unreliable at small $n$, motivating
warning-stratified evaluation as standard practice for safety-sensitive
NLP, not an optional addition.

Taken together with the classifier results, the preliminary picture
matches our hypothesis: rule-based substitution is safe but limited to
the $\sim\!43\%$ of edits that are Substitution-type, while unguided LLM
prompting performs the full range of edits but occasionally sacrifices
safety-critical content---exactly the failure mode our operation-aware
router is designed to prevent.


\section{Conclusion and Next Steps}

Our preliminary results support the diagnostic framing of this paper:
random-sample evaluation alone suggests all baselines preserve safety
equally, but a warning-stratified stress test reveals unguided LLM
prompting uniquely drops safety-critical content while achieving the
highest readability--the trade-off our operation-aware system is
designed to break. Our TF-IDF classifier, despite using pseudo-labels,
performs within the literature's reported macro-F1 range for this task
\citep{ondov2025lessons}, with the strongest confusion between
Substitution and Explanation rather than Explanation and Generalization
as initially expected, motivating contextual models next.

For the final paper we will add BERT-base and BioBERT classifiers,
evaluate the full operation-aware generation pipeline against all three
baselines on the complete test set, extend the safety metrics to entities
and numerical expressions, and integrate full QuickUMLS lookup. If the
operation-aware system Pareto-dominates the baselines on the
readability--safety plane, that is the finding the final paper reports.

\begin{table*}[t]
  \centering
  \small
  \begin{tabular}{lccccc}
    \toprule
    \textbf{System} & \textbf{SARI}$\uparrow$ & \textbf{FKGL}$\downarrow$ &
    \textbf{Compr.} & \textbf{Warn.\ Pres.} & \textbf{Warn.\ Pres.}\textsuperscript{$\dagger$} \\
     & & & & \textbf{(random, $n$=50)} & \textbf{(stress test, $n$=30)} \\
    \midrule
    No simplification    & 16.82 & 15.31 & 1.00  & 1.000 & 1.000 \\
    Rule-based CHV        & 18.66 & 15.17 & 1.02  & 1.000 & 0.994 \\
    Direct LLM (FLAN-T5)  & 22.28 & 14.89 & 0.945 & 1.000 & 0.950 \\
    \bottomrule
  \end{tabular}
  \caption{Preliminary results on PLABA validation data. SARI/FKGL/compression
    from a random 50-sentence sample; warning preservation reported for both
    the same random sample and a warning-bearing stress test ($n=30$).
    \textsuperscript{$\dagger$}The stress test isolates the safety failure
    invisible to random sampling---see Section~\ref{sec:metrics}.}
  \label{tab:prelim_results}
\end{table*}

\section{Related Work}

\paragraph{Datasets and prior PLABA systems.}
PLABA \citep{attal2023plaba} pairs biomedical abstracts with expert
plain-language rewrites at the sentence level, letting us pseudo-label
operations rather than treat simplification as opaque sequence-to-sequence
generation. \citet{knappich2023boschai}'s PLABA-2023-winning BoschAI
system finds that fine-tuned models copy the input and edit
little---directly motivating our diagnostic framing, since a model that
barely edits also barely distinguishes a safe substitution from a
critical warning. \citet{li2024beemanc} investigate LLMs with
controllable attributes for the same track, including the BART-based and
Lay-SciFive systems we plan as neural baselines in the full paper.

\paragraph{TSAR 2025 and the evaluation gap.}
The winning TSAR 2025 submission \citep{miyata2025ehimenlp} generates
candidates via 28 model-prompt combinations and reranks by semantic
similarity; their own ablation shows the gain comes from diversity and
selection, not from understanding \emph{what kind} of edit a span needs.
The shared task's own findings \citep{alvamanchego2025tsarfindings} show
that across 48 submissions from 20 teams, none classified complexity
type before generating output, and the official metrics (CEFR compliance,
MeaningBERT similarity) cannot detect a fluent output that quietly drops
a warning or dosage---the organizers explicitly call for
expert-in-the-loop validation as future work. Our preservation metrics
answer this call by making safety a first-class, automatically measurable
criterion rather than deferring it to expert review.

\paragraph{Jargon detection and style robustness.}
\citet{papandreou2025jargon} show jargon detectors trained on MedReadMe
transfer poorly to PLABA (33.71\% entity F1), since the datasets define
jargon under different annotation objectives---supporting our decision to
build detectors natively on PLABA. Their jargon-aware LLM prompts also
introduce factual errors a lookup table structurally cannot produce,
motivating our CHV-based Substitution pathway over jargon-surfacing
prompts. \citet{rahman2025spqa} find that stylistic perturbations
significantly degrade LLM QA correctness and completeness even when
fluency stays stable, external evidence that surface-level metrics can
miss real content degradation---the same gap our safety metrics target.

\paragraph{Prior operation-aware simplification.}
\citet{cripwell2022operation} propose controllable simplification via
operation classification as a predict-then-execute process, the same
structural intuition we extend to the biomedical span level.
\citet{devaraj2021medical} identify safety-critical information loss as a
recurring failure in paragraph-level medical simplification, directly
motivating our preservation metrics. \citet{xia2025jebs} introduce JEBS
on the same five-operation PLABA taxonomy, arguing that separating
identification, classification, and generation enables more targeted
evaluation---the diagnose-before-simplify principle underlying our
pipeline.

\paragraph{Evaluation.}
SARI \citep{xu2016sari} rewards correct add/keep/delete operations against
multiple references; BERTScore \citep{zhang2020bertscore} adds contextual
semantic similarity. Neither captures whether a warning or dosage
survived, consistent with reports that SARI and FKGL correlate only
moderately with human judgment \citep{alvamanchego2021unsuitability}---the
gap our preservation metrics are built to fill.

\section*{Limitations}

Our span-level operation labels are pseudo-labels derived from length and
edit-distance heuristics, not gold annotations; sentences that
simultaneously add explanation and remove specificity can be mislabeled
when length changes cancel. This is consistent with the broader
difficulty of the task: \citet{ondov2025lessons} report human
inter-annotator agreement of only 0.46 F1 on operation classification,
suggesting some divergence between our pseudo-labels and true operation
categories is an expected property of the task rather than a defect
specific to our heuristic. Our complexity detectors rely on handcrafted
patterns and therefore report a lower bound on true complexity, and the
UMLS matcher currently uses a morphological fallback rather than a full
UMLS lookup. The warning-preservation metric relies on a curated
paraphrase map, which encodes annotator assumptions; we plan to replace
it with embedding-based semantic matching. Finally, preliminary results
use a small validation sample and a single classifier rung; the
contextual models and full-corpus evaluation remain future work.

\bibliography{custom}

\end{document}
